% =============================================================================
% main.tex — A Survey on Non-Line-of-Sight Optical Camera Communication
% Companion to refs.bib (36 references)
% =============================================================================
\documentclass[10pt,a4paper]{article}

\usepackage[utf8]{inputenc}
\usepackage[T1]{fontenc}
\usepackage{cite}
\usepackage{amsmath,amssymb}
\usepackage{graphicx}
\usepackage{booktabs}
\usepackage{multirow}
\usepackage{url}
\usepackage{hyperref}
\usepackage{xcolor}
\usepackage{geometry}
\usepackage{array}
\usepackage{tabularx}

\geometry{margin=2.2cm}

\hypersetup{
  colorlinks=true,
  linkcolor=blue!60!black,
  citecolor=teal!70!black,
  urlcolor=blue!60!black
}

% Convenience macros
\newcommand{\rp}[1]{\textbf{RP#1}}
\newcommand{\etc}{\textit{etc.}\ }
\newcommand{\eg}{\textit{e.g.}\ }
\newcommand{\ie}{\textit{i.e.}\ }

\title{Non-Line-of-Sight Optical Camera Communication: \\
       A Survey of Research Frontiers and Technical Details}

\author{
  NLOS-OCC Research Group
}

\date{September 2026}

% =============================================================================
\begin{document}
\maketitle

% -----------------------------------------------------------------------------
\begin{abstract}
Non-Line-of-Sight Optical Camera Communication (NLOS-OCC) extends
conventional optical camera communication beyond direct visibility by
exploiting diffuse reflections, partial occlusions, and ambient-light
scattering to recover data from light-emitting diode (LED) signals that
never reach the camera along a free-space path. The central technical
contradiction is that LEDs are modulated in the MHz range while consumer
cameras sample at 30--60 frames per second, so a single frame cannot
resolve individual symbols directly. This survey organizes the rapidly
expanding NLOS-OCC literature into a five-layer research framework
comprising twelve research points ($\rp{1}$--$\rp{12}$): physical-layer
channel modelling, modulation and coding, signal recovery, hardware and
receiver front-ends, networking, and applications. For each layer we
distil the key technical mechanisms---Barker-code pilots, differential
dual-bit processing, rolling-shutter stripe extraction, U-Net bit
segmentation, generative-adversarial background/stripe separation,
event-camera asynchronous detection, reconfigurable intelligent surfaces
(RIS), MIMO and NOMA schemes, and learning-based smartphone OCC
architectures. We further survey four emerging directions
($\rp{13}$--$\rp{16}$): UAV-mounted OCC, underwater and cross-medium
OCC, semantic optical communication toward AI-native 6G, and
retroreflective optical integrated sensing and communication (RO-ISAC).
We further examine the recently approved IEEE~802.15.7a-2024 amendment
whose PHY~VII defines seven OCC modes targeting 200\,m range and
350\,km/h mobility. We close with a prioritized roadmap that traces the
dependency chain from channel characterization through signal recovery
and networking to end-user applications, identifying near-, mid-, and
long-term research priorities.
\end{abstract}

\noindent\textbf{Index Terms}---Optical camera communication (OCC),
non-line-of-sight (NLOS), visible light communication (VLC), rolling
shutter, deep learning, event camera, MIMO, RIS, IEEE 802.15.7a,
vehicular communication, Internet of Things, UAV, underwater OCC,
semantic communication, integrated sensing and communication (ISAC).

% =============================================================================
\section{Introduction}
\label{sec:intro}

Optical Camera Communication (OCC) is a branch of visible light
communication (VLC) in which the receiver is an image sensor rather
than a photodiode~\cite{ieee8021572018}. By spatially resolving light
sources, a camera can discriminate among multiple transmitters and
reject ambient interference in a way that single-element receivers
cannot. The cost of this spatial selectivity is a severe temporal
sampling limitation: a commodity camera captures 30--60 frames per
second, while the LED driving signal may be modulated at MHz
frequencies. The rolling-shutter effect, whereby successive rows of the
sensor are exposed at slightly different times, partially alleviates
this by spreading a time-varying signal across the image rows, producing
characteristic light-and-dark \emph{stripes}; but it cannot match the
throughput of photodiode-based VLC.

The challenge becomes even more acute when the line-of-sight path
between the LED and the camera is blocked. In Non-Line-of-Sight OCC
(NLOS-OCC), the camera must recover data from diffuse reflections off
walls, the ground, or other surfaces, or from partially occluded sources
seen through glass or foliage. The received optical power is drastically
reduced, the signal is smeared by the bidirectional reflectance
distribution function (BRDF) of the reflecting surface, and the
background illumination is heterogeneous. Yet NLOS operation is
essential for practical deployment: in vehicular scenarios vehicles are
frequently obscured by other traffic; in indoor scenarios furniture and
occupants block direct paths; and in smart-city deployments street
furniture creates persistent occlusions.

This survey organizes the NLOS-OCC literature into a five-layer research
framework (Fig.~\ref{fig:framework}) and twelve core research points
($\rp{1}$--$\rp{12}$), summarized in Table~\ref{tab:rps}, plus four
emerging directions ($\rp{13}$--$\rp{16}$) surveyed in
Section~\ref{sec:emerging}. The framework
mirrors the protocol stack: the physical layer ($\rp{1}$--$\rp{3}$)
covers channel modelling, modulation, and coding; the signal-recovery
layer ($\rp{4}$--$\rp{6}$) covers noise suppression, deep-learning
extraction, and event-camera detection; the hardware layer
($\rp{7}$--$\rp{8}$) covers smartphone receivers and RIS-enhanced links;
the network layer ($\rp{9}$--$\rp{10}$) covers MIMO and heterogeneous
fusion; the application layer ($\rp{11}$--$\rp{12}$) covers
V2V/V2I/IoT systems; and the emerging layer ($\rp{13}$--$\rp{16}$)
covers UAV, underwater, semantic, and ISAC frontiers.

\begin{table*}[t]
\centering
\caption{The NLOS-OCC research framework: sixteen research points across six layers.}
\label{tab:rps}
\small
\begin{tabularx}{\textwidth}{l l X}
\toprule
\textbf{Layer} & \textbf{RP} & \textbf{Focus} \\
\midrule
Physical ($\rp{1}$--$\rp{3}$) & $\rp{1}$ & NLOS-OCC channel modelling (BRDF, ray tracing, geometry) \\
 & $\rp{2}$ & Modulation schemes (OOK, PWM, UFSOOK, CSK, CSS, N-pulse, adaptive) \\
 & $\rp{3}$ & Channel coding (Raptor/RaptorQ fountain codes, RS, FEC) \\
\midrule
Signal Recovery ($\rp{4}$--$\rp{6}$) & $\rp{4}$ & Noise suppression and adaptive thresholding \\
 & $\rp{5}$ & Deep-learning signal recovery (U-Net, CNN, GAN, autoencoder) \\
 & $\rp{6}$ & Event-camera / neuromorphic asynchronous detection \\
\midrule
Hardware ($\rp{7}$--$\rp{8}$) & $\rp{7}$ & Smartphone OCC (Pro mode, RAW domain, deep-learning demodulation) \\
 & $\rp{8}$ & RIS / ORIS for NLOS coverage enhancement \\
\midrule
Network ($\rp{9}$--$\rp{10}$) & $\rp{9}$ & MIMO schemes (WDM-MIMO, SDM, 2$\times$N, MIMO-NOMA) \\
 & $\rp{10}$ & Heterogeneous fusion (OCC/BLE, OCC/WiFi, 6G, IEEE 802.15.7a) \\
\midrule
Applications ($\rp{11}$--$\rp{12}$) & $\rp{11}$ & V2V / V2I / ITS vehicular systems \\
 & $\rp{12}$ & Indoor VLP, IoT, ISAC, smart-city deployment \\
\midrule
Emerging ($\rp{13}$--$\rp{16}$) & $\rp{13}$ & UAV-mounted OCC (drone receiver, 360$^\circ$ FOV) \\
 & $\rp{14}$ & Underwater and cross-medium (air-water) OCC \\
 & $\rp{15}$ & Semantic optical communication, AI-native 6G integration \\
 & $\rp{16}$ & Retroreflective optical ISAC (RO-ISAC) \\
\bottomrule
\end{tabularx}
\end{table*}

\begin{figure}[t]
  \centering
  \fbox{\parbox{0.92\columnwidth}{\centering\small
  \textbf{Emerging} ($\rp{13}$--$\rp{16}$): UAV, underwater, semantic 6G, RO-ISAC\\
  $\updownarrow$\\
  \textbf{Application} ($\rp{11}$,$\rp{12}$): V2V/V2I, IoT, ISAC\\
  $\updownarrow$\\
  \textbf{Network} ($\rp{9}$,$\rp{10}$): MIMO, heterogeneous fusion, 802.15.7a\\
  $\updownarrow$\\
  \textbf{Hardware} ($\rp{7}$,$\rp{8}$): smartphone OCC, RIS\\
  $\updownarrow$\\
  \textbf{Signal Recovery} ($\rp{4}$--$\rp{6}$): denoising, deep learning, event cameras\\
  $\updownarrow$\\
  \textbf{Physical} ($\rp{1}$--$\rp{3}$): channel, modulation, coding}}
  \caption{Five-layer NLOS-OCC research framework. Information flows
  upward; constraints propagate downward. The dependency chain
  $\rp{1}\rightarrow\rp{4}/\rp{6}\rightarrow\rp{9}\rightarrow\rp{11}$
  defines the critical research path.}
  \label{fig:framework}
\end{figure}

The remainder of the paper follows the framework layer by layer
(Sections~\ref{sec:physical}--\ref{sec:applications}), surveys four
emerging research directions (Section~\ref{sec:emerging}), examines the
IEEE~802.15.7a-2024 standard (Section~\ref{sec:standard}), and closes
with a prioritized roadmap and conclusions
(Sections~\ref{sec:roadmap}--\ref{sec:conclusion}).

% =============================================================================
\section{Physical Layer: Channel, Modulation, and Coding}
\label{sec:physical}

\subsection{NLOS-OCC Channel Modelling ($\rp{1}$)}
\label{sec:rp1}

The NLOS-OCC channel departs fundamentally from the LOS VLC channel.
Instead of a direct irradiance link, the received optical power is the
integral of the LED radiance reflected by one or more surfaces whose
BRDF may be Lambertian, specular, or mixed. Three modelling approaches
dominate the literature.

\paragraph{Analytical BRDF models.}
A Lambertian reflector re-emits incident light uniformly in all
directions with radiance proportional to the cosine of the exit angle.
Real surfaces---asphalt, painted walls, foliage---deviate from this
ideal, and the deviation directly affects the spatial distribution of
received power across the image sensor. The pilot-aided NLOS-OCC work of
Chen \emph{et al.}~\cite{chen2021pilot} measured a 2.5\,m NLOS link under
heterogeneous background light, demonstrating that the pilot-based
signal-logging and estimation (PSLE) scheme combined with enhanced
visual compensation (EVC) can recover data where unaided thresholding
fails entirely. The experiment confirms that the reflected-signal
contrast ratio is typically 5--15\,dB below the LOS case, placing a hard
floor on the achievable stripe signal-to-noise ratio.

\paragraph{Ray-tracing simulation.}
For vehicular MIMO channels, Liu and Dang~\cite{liu2023vehicular}
proposed a three-dimensional irregular-shaped geometry-based stochastic
model (3D IS-GBSM) that captures the birth-death process of scattering
clusters as vehicles move. The model parameters are derived from
OpticStudio ray tracing and validated against measured
power-delay profiles. The same ray-tracing approach underlies the
MIMO V2V-VLC study of Yahia \emph{et al.}~\cite{yahia2022mimo}, which
simulates dual-transmitter headlight configurations and weather-dependent
attenuation (clear, rain, fog). A key finding is that MIMO spatial
multiplexing gain collapses rapidly as fog visibility drops below
100\,m, motivating the adaptive and coding strategies discussed in
Section~\ref{sec:rp3}.

\paragraph{Pilot-based long-distance NLOS.}
Dong \emph{et al.}~\cite{dong2023barker} achieved a 5\,m NLOS-OCC link
using a Barker-code pilot and a differential dual-bit processing (DDBP)
algorithm. The Barker code's autocorrelation property lets the receiver
locate the pilot in the stripe sequence even when the reflected signal
is deeply buried in ambient light; DDBP then cancels the slowly varying
background by subtracting consecutive bit windows. The 7\% forward-error-
correction (FEC) threshold was met at 5\,m, extending the practical
NLOS range beyond earlier pilot-free demonstrations.

\subsection{Modulation Schemes ($\rp{2}$)}
\label{sec:rp2}

OCC modulation must respect the camera's sampling constraints. The
principal families are:

\begin{itemize}
  \item \textbf{OOK / PWM.} On-off keying and pulse-width modulation are
    the simplest schemes. With rolling shutter, each symbol produces a
    band of bright or dark rows. The symbol rate is limited to roughly
    $f_s / 2$ where $f_s$ is the effective row-sampling frequency
    (typically a few kHz).
  \item \textbf{UFSOOK.} Undersampled Frequency Shift OOK modulates at a
    frequency that aliases to a distinguishable low frequency under the
    camera's sampling, enabling per-frame demodulation without relying
    on rolling shutter. The MIMO extension~\cite{mimoufsook2013} maps
    different frequencies to spatially separated LEDs, achieving a
    rudimentary spatial-multiplexing gain.
  \item \textbf{CSK / Color barcode.} Color Shift Keying encodes data in
    the colour coordinate of the emitted light. COBRA~\cite{cobra2012}
    demonstrated real-time colour-barcode streaming to smartphones, and
    the approach extends naturally to multi-colour LED arrays. CSK's
    chromaticity constraint (perceptual flicker avoidance) limits the
    practical constellation size.
  \item \textbf{Spatial modulation / PixNet.} PixNet~\cite{pixnet2010}
    treats an LED array as a 2D spatial modulator, encoding blocks of
    bits in the on/off pattern of individual elements. This is the
    screen-camera analog of MIMO spatial modulation and achieves the
    highest raw rates among OCC schemes, but requires a well-resolved
    array image---difficult in NLOS conditions.
  \item \textbf{Adaptive modulation.} Shirvani \emph{et al.}
    \cite{shirvani2021adaptive} investigated adaptive thresholding for
    V2V OCC where a mobile-phone camera receives LED taillight signals.
    The threshold is adjusted per-frame based on the measured stripe
    intensity statistics, accommodating the rapid variation in channel
    gain as the vehicles' relative position changes.
\end{itemize}

For NLOS links specifically, the modulation choice is constrained by the
low contrast of the reflected signal. Pure OOK with a fixed threshold
fails because the reflected intensity varies with surface geometry;
adaptive and differential schemes (DDBP, PSLE) are essential.

\subsection{Channel Coding ($\rp{3}$)}
\label{sec:rp3}

Because the NLOS-OCC channel suffers deep fades from surface
roughness, transient occlusions, and frame drops, rateless and fountain
codes are particularly attractive: they generate as many encoded
symbols as needed until the receiver acknowledges successful decoding,
gracefully adapting to the unknown channel quality.

Raptor codes were applied to free-space optical communication by Lu
\emph{et al.}~\cite{lu2016raptor}, and RaptorQ---a more robust variant
with near-MDS (minimum distance separable) performance---was studied
under atmospheric turbulence by Pernice \emph{et al.}
\cite{pernice2016raptorq}. RaptorQ's advantage in the OCC context is
its ability to recover the original block from any sufficiently large
subset of received symbols, tolerating the bursty erasures that arise
when a vehicle passes behind an obstruction and the camera drops several
consecutive frames.

The IEEE 802.15.7a-2024 amendment (Section~\ref{sec:standard}) adopts
Reed-Solomon (RS) codes in its RS-OFDM mode, combined with convolutional
or LDPC inner codes, providing a standardized alternative to the
research-grade RaptorQ approach.

% =============================================================================
\section{Signal Recovery Layer}
\label{sec:signal}

\subsection{Noise Suppression and Adaptive Thresholding ($\rp{4}$)}
\label{sec:rp4}

In NLOS-OCC the signal of interest---the faint stripe pattern from a
reflected LED---is superimposed on a spatially varying background
comprising ambient illumination, other light sources, and the reflecting
surface's own texture. Classical thresholding assumes a stationary
background; in NLOS this assumption fails.

The pilot-aided approaches ($\rp{1}$) already incorporate background
cancellation (DDBP, EVC). Beyond these, adaptive per-stripe thresholding
\cite{shirvani2021adaptive} estimates the local mean and variance from
the stripe's own row profile and sets a threshold at the midpoint,
yielding robust demodulation under vehicle motion. For quasi-NLOS
screen-camera links through glass, Wang \emph{et al.}
\cite{screencamera2015} demonstrated that the glass reflection creates a
secondary image that can be separated by disparity cues, recovering
communication behind the glass. For long-distance outdoor links where
atmospheric turbulence causes intensity fluctuations, Li \emph{et al.}
\cite{li2026manchester} proposed a Manchester-coding-based demodulation
that exploits the guaranteed level transition in each symbol period,
making the decoder robust to slow scintillation that would defeat a
fixed threshold.

\subsection{Deep-Learning Signal Recovery ($\rp{5}$)}
\label{sec:rp5}

The last several years have seen a rapid migration from hand-crafted
signal processing to learned pipelines. The principal architectures are:

\paragraph{U-Net bit segmentation.}
A U-Net semantic-segmentation network can be trained to produce a
pixel-level binary mask classifying each pixel as ``bright stripe'' or
``dark stripe'' \cite{unet2025nlos}. The encoder--decoder structure with
skip connections preserves the fine spatial detail of the stripe edges,
which a plain CNN classifier tends to blur. Once the mask is obtained,
the bit sequence follows by simple row-major reading. This approach is
particularly effective for NLOS reflections where the stripe boundaries
are diffuse and the contrast is low.

\paragraph{GAN-based background/stripe separation.}
When the background is dynamic---moving vehicles, changing illumination,
flickering displays---a fixed or slowly adaptive threshold cannot
separate signal from background. A generative adversarial network can be
trained to decompose the captured image into a stripe component and a
background component, the discriminator enforcing that the generated
background is indistinguishable from real background-only frames
\cite{ganocc2026}. The separated stripe image is then demodulated with
conventional or learned decoders. The GAN approach has been shown to
reduce bit-error rate under motion and partial occlusion, achieving
usable NLOS communication at 2.5\,m with $10^{-4}$ BER in the
remote-control demonstration of Ahmed \emph{et al.} (DOI
10.1109/JPHOT.2023.3250455).

\paragraph{Two-stage and end-to-end DNN.}
A two-stage pipeline---LED detection followed by stripe
classification---reduces the search space before demodulation, improving
both accuracy and latency. End-to-end autoencoders that map bits to LED
drive signals and images back to bits have been explored for LOS OCC and
are a promising direction for NLOS, though training-data scarcity for
the NLOS channel remains an obstacle. For multi-level Colour Shift
Keying (CSK), Furukawa \emph{et al.} \cite{furukawa2025cskdl} showed
that a deep-learning equalizer can compensate the constellation
distortion caused by hue shifts of the light source and the image
sensor's internal colour processing, recovering the nominal CSK
constellation points and improving demodulation accuracy---an important
step toward high-spectral-efficiency OCC beyond binary modulation.

\subsection{Event Cameras and Neuromorphic Detection ($\rp{6}$)}
\label{sec:rp6}

The most architecturally radical approach to the sampling bottleneck is
to abandon the frame-based camera entirely. An event camera (Dynamic
Vision Sensor, DVS) fires an asynchronous pixel-level event whenever the
log-intensity at that pixel changes by more than a threshold. Because
the LED's high-frequency modulation produces continuous intensity
changes, the event camera effectively samples the signal at
$\sim$20\,kHz---three orders of magnitude above frame-based cameras.

The neuromorphic NLOS-OWC work \cite{neuromorphic2025nlos} (arXiv
2503.11226) is the first demonstration of event-camera-based NLOS optical
wireless communication. Key results include an effective sampling rate
of $\sim$20\,kHz and a dynamic range of 120\,dB, the latter being
critical for NLOS where the reflected signal may span several orders of
magnitude as the reflecting surface geometry changes. The motion-aware
event-camera OCC study \cite{event2024motion} further shows that the
asynchronous detection is inherently tolerant to transmitter and receiver
motion, a property that is highly desirable for vehicular deployment.

The trade-off is that event cameras produce sparse, asynchronous data
streams that require fundamentally different processing pipelines
(spiking neural networks, event-to-frame conversion, or direct
event-based decoding) compared to the image-centric pipelines of
$\rp{4}$--$\rp{5}$.

% =============================================================================
\section{Hardware Layer}
\label{sec:hardware}

\subsection{Smartphone OCC: Pro Mode, RAW Domain, and Learning-Based Demodulation ($\rp{7}$)}
\label{sec:rp7}

The ultimate deployment target for much OCC research is the commodity
smartphone, whose camera is ubiquitous but whose image signal processor
(ISP) is optimized for photography, not communication. The ISP applies
auto-exposure, auto-white-balance, demosaicing, denoising, and
JPEG/H.265 compression---each of which can destroy the faint stripe
pattern that carries data.

\paragraph{WinkLink / Xiao \emph{et al.}}
The MobiSys~2024 work of Xiao \emph{et al.}~\cite{winklink2024} tackles
the hardest practical case: optical camera communication when the LED
emitter is \emph{unseen} (outside the camera's field of view or occluded
by complex, non-uniform backgrounds). The system employs a two-stage
deep neural network. The first stage performs context-aware
demodulation: it learns to separate the faint data-bearing stripe
pattern from arbitrary background clutter (posters, textured walls,
mixed illumination) that would defeat any fixed-threshold detector. The
second stage reconstructs the transmitted bits from the cleaned signal.
Crucially, the network is trained entirely on \emph{synthesized} data
---rendered stripes composited onto photographed backgrounds---yet
generalizes to the real world: it was validated in 12 distinct
environments and tolerates device motion during capture. Quantitatively,
the system operates at $-20$\,dB SNR, delivers a $5.8$\,dB SNR gain over
the best baseline, and extends the effective communication distance by
$5.5\times$ (reaching 11\,m with a 10\,W LED). This data-driven,
learning-centric approach represents a paradigm shift from the
hand-crafted signal-processing pipelines that dominate earlier OCC
literature, and it directly addresses the NLOS ``complex background''
problem that is the central challenge of this survey.

\paragraph{Passive and ambient-light approaches.}
PassiveVLC~\cite{passivevlc2017} (MobiCom 2017) takes a complementary
approach: rather than driving an LED, it modulates the reflectance of a
passive tag illuminated by ambient light, enabling backscatter
communication with zero transmit energy. LightSync~\cite{lightsync2013}
(MobiCom 2013) uses ambient light fluctuations for passive inter-device
clock synchronization, a building block for any OCC system that must
align transmitter and receiver timing without a shared reference.

\paragraph{Pro mode and 4CASK-OCC.}
Modern smartphone camera APIs (Android Camera2, iOS AVFoundation) expose
manual exposure, focus, and ISO controls. Exploiting these ``Pro mode''
settings---short exposure, fixed white balance, high ISO---substantially
improves stripe visibility. The 4CASK-OOC scheme (four-colour CSK)
extends the colour-barcode idea to the smartphone's Bayer pattern,
mapping symbols to quadrants of the CIE colour space to maximize
perceptual-flicker-free throughput.

\subsection{RIS / ORIS for NLOS Coverage ($\rp{8}$)}
\label{sec:rp8}

When no reflecting surface is available, or when the natural reflection
is too weak, a Reconfigurable Intelligent Surface (RIS) can be deployed
to create an artificial NLOS path. An RIS is a planar array of
sub-wavelength elements whose reflection phase can be controlled
electronically, allowing the surface to steer an incident beam toward
the receiver even when the transmitter and receiver are not mutually
visible.

The RIS-ISAC study \cite{risisac2025v2x} formulates the joint
beamforming design for an RIS-assisted integrated sensing and
communication (ISAC) V2X system, minimizing the Cram\'er-Rao bound of
the sensing estimate subject to SINR constraints on the communication
links. Reported coverage enhancements are 15--20\,dB in the NLoS region,
transforming an unusable blocked link into one with margin to spare. For
OCC specifically, an outdoor RIS (ORIS) can redirect a vehicle's
headlight modulation toward a roadside camera that would otherwise be in
shadow, bridging the coverage gap at intersections and in urban canyons.

% =============================================================================
\section{Network Layer}
\label{sec:network}

\subsection{MIMO Schemes ($\rp{9}$)}
\label{sec:rp9}

Multiple-input multiple-output (MIMO) techniques increase throughput and
reliability by exploiting the spatial dimension. In OCC, the camera's
spatial resolution is a natural enabler: each LED in an array can carry
an independent data stream, and the image sensor separates them by
location.

\paragraph{WDM-MIMO and SDM.}
Wavelength-division multiplexing (WDM) assigns different colours
(R/G/B) to different streams; space-division multiplexing (SDM) assigns
different physical LEDs. Both are limited by the camera's colour
separation and spatial resolution, respectively.

\paragraph{2$\times$N MIMO and ray tracing.}
Yahia \emph{et al.}~\cite{yahia2022mimo} simulated a 2$\times$2 MIMO
V2V-VLC system using dual-transmitter headlight configurations and
OpticStudio ray tracing, quantifying the channel cross-talk under
clear, rain, and fog conditions. The results show that zero-forcing (ZF)
spatial demultiplexing is adequate in clear weather but degrades as
weather reduces the rank of the channel matrix.

\paragraph{MIMO-NOMA.}
The combination of MIMO with non-orthogonal multiple access (NOMA) is a
powerful approach for V2V-VLC. In \cite{mimonoma2025v2v}, a 2$\times$2
MIMO system uses power-domain NOMA with successive interference
cancellation (SIC), DCO-OFDM for bandwidth efficiency, and a novel
normalized gain difference power allocation (NGDPA) scheme. The reported
throughput improvement is 20\% over OFDMA. A related beamforming study
\cite{vlcNoma2023beamforming} applies NOMA with SIC and beamforming in
the V2V context, showing further gains when the beamformer aligns the
power-domain layers with the channel's spatial modes.

\paragraph{MIMO survey.}
Sejan \emph{et al.}~\cite{sejan2023survey} provide a comprehensive
survey of MIMO-VLC covering the full taxonomy (multiplexing, diversity,
multi-user), machine-learning methods for channel estimation and
demodulation, and future trends toward massive MIMO and intelligent
reflecting surfaces.

\subsection{Heterogeneous Fusion ($\rp{10}$)}
\label{sec:rp10}

Pure OCC links are brittle: they fail in direct sunlight, in heavy fog,
and when the camera's view is fully obstructed. Heterogeneous fusion
combines OCC with a complementary radio technology to provide a robust
overlay.

\paragraph{OCC/VLC + RF.}
Okasha \emph{et al.}~\cite{okasha2022hybrid} propose a hybrid car-to-car
MIMO VLC+RF system in which each vehicle alternates between a 100\,m VLC
hop and a 200\,m RF hop, supporting up to 5 relay hops. The VLC hop
achieves four times the throughput of the RF hop when the optical link
is available, and the RF hop maintains connectivity when it is not.

\paragraph{OCC/BLE and OCC/WiFi.}
In indoor IoT scenarios, fusing OCC with Bluetooth Low Energy (BLE) or
WiFi provides seamless handover: OCC delivers high-rate downlink data to
a smartphone camera while BLE/WiFi handles the low-rate uplink and
control channel. This split is natural because smartphones lack an
uplink light source suitable for communication.

\paragraph{6G and IEEE 802.15.7a.}
The IEEE 802.15.7a-2024 amendment (Section~\ref{sec:standard})
formalizes OCC as a PHY within the short-range optical wireless
standard, providing the interoperability layer needed for heterogeneous
6G networks. The standard's relaying mechanism enables an OCC link to
hand off to a neighbouring OCC or RF node, implementing the fusion
architectures above in a standardized way.

% =============================================================================
\section{Application Layer}
\label{sec:applications}

\subsection{Vehicular Communication: V2V, V2I, and ITS ($\rp{11}$)}
\label{sec:rp11}

Vehicular visible light communication is the flagship application of
OCC. Yamazato \emph{et al.}~\cite{yamazato2014vehicular} laid the
foundation by demonstrating V2V and V2I links using vehicle LED
headlights and taillights, highlighting the road-surface reflection
channel that naturally arises in urban driving. Aly \emph{et al.}
\cite{aly2022vehicular} improved robustness with selection combining
receiver diversity, selecting the best of multiple photodetector
branches to combat the rapid fading of the vehicular channel.

For camera-based receivers, the V2V OCC system of Shirvani \emph{et al.}
\cite{shirvani2021adaptive} uses a mobile-phone camera pointed at the
preceding vehicle's taillights, with adaptive thresholding to handle the
changing channel. The MIMO-NOMA V2V-VLC system
\cite{mimonoma2025v2v} pushes raw rates higher with
2$\times$2 spatial multiplexing and power-domain superposition, while
the hybrid VLC+RF architecture \cite{okasha2022hybrid} extends range
beyond a single hop. The RIS-ISAC V2X design \cite{risisac2025v2x} adds
intelligent reflecting surfaces at intersections to maintain coverage
when buildings block the direct path.

\subsection{Indoor VLP, IoT, and ISAC ($\rp{12}$)}
\label{sec:rp12}

Beyond vehicles, NLOS-OCC serves indoor visible light positioning (VLP),
Internet of Things (IoT) sensor connectivity, and integrated sensing and
communication (ISAC). Indoor VLP leverages the camera's angular
resolution to estimate the position of an LED beacon from its image,
achieving sub-metre accuracy even when the beacon is seen only as a
diffuse reflection on a wall. IoT applications use low-rate OCC to read
sensor values from LED indicators on appliances, with the smartphone
camera as the universal reader. ISAC combines the communication function
with a sensing function (presence detection, gesture recognition) using
the same LED illumination, an approach that aligns with the RIS-ISAC
framework of \cite{risisac2025v2x} and the AI-PHY provisions of the
IEEE 802.15.7a standard.

% =============================================================================
\section{Emerging Research Directions}
\label{sec:emerging}

Beyond the twelve research points of the core framework, four emerging
directions are reshaping the NLOS-OCC landscape. These directions extend
the framework upward (toward 6G semantic integration), outward (to UAV
and underwater environments), and inward (to joint sensing-communication
waveforms). We label them $\rp{13}$--$\rp{16}$ to indicate their
position as frontier extensions of the five-layer stack.

\subsection{UAV-Mounted OCC ($\rp{13}$)}
\label{sec:rp13}

Unmanned aerial vehicles (UAVs) equipped with cameras or LED payloads
open a new deployment axis for OCC. Chang \emph{et al.}
\cite{chang2026uavocc} demonstrated the first wide-field-of-view (FOV)
bidirectional UAV-assisted OWC system using the under-sampled phase
shift on-off keying (UPS-OOK) format and a light-diffusing-fibre
transmitter that radiates over 360$^\circ$, enabling a hovering UAV to
serve ground receivers at any azimuth. The long-short-term-memory neural
network (LSTM-NN) decodes the rolling-shutter pattern from arbitrary
viewing angles, removing the pointing constraint that burdens
conventional OCC.

For agricultural IoT, Kondo \emph{et al.} \cite{kondo2024farmocc}
proposed a drone-based OCC sensor network in which LED panels on ground
sensor nodes transmit soil-moisture and temperature data to a
drone-mounted camera. A trajectory-control algorithm reduces total
flight time by 30\,\% versus a shortest-path baseline, and YOLOv5-based
LED detection with simple OOK modulation achieves BER $<10^{-3}$ in a
real-world leaf-mustard farm trial. This paradigm---low-cost,
radio-silent, infrastructure-free---is attractive for underdeveloped and
radio-quiet zones.

\subsection{Underwater and Cross-Medium OCC ($\rp{14}$)}
\label{sec:rp14}

Underwater OCC (UWOCC) inherits the NLOS challenge in an extreme form:
the water column introduces wavelength-dependent absorption, particulate
scattering, and turbulence that jointly attenuate and smear the optical
signal. Hamidnejad and Gholami \cite{hamidnejad2023uwmimoocc} developed
a comprehensive underwater MIMO OCC channel model that captures these
effects and showed that spatial modulation---exploiting the camera's
ability to resolve multiple LED transmitters---retains its throughput
advantage even in turbid water.

At the air-water interface, Tsai \emph{et al.} \cite{tsai2024uavwatertoair}
combined the UAV and underwater themes: a UAV-mounted camera receives
OCC signals transmitted from an underwater source, crossing the wavy,
cloudy water surface. An LSTM-NN mitigates the surface-fluctuation-induced
link outage, demonstrating that cross-medium NLOS OCC---from a submerged
sensor to an aerial camera---is feasible. This opens paths for
autonomous-underwater-vehicle (AUV) telemetry retrieval without
recovering the vehicle to the surface.

\subsection{Semantic Optical Communication and AI-Native 6G ($\rp{15}$)}
\label{sec:rp15}

Semantic communication shifts the fidelity criterion from bit-level
reconstruction to task-relevant meaning, trading redundant detail for
robustness. For underwater image transmission, Nennouche \emph{et al.}
\cite{nennouche2025semanticuwoc} proposed a variational-autoencoder
(VAE) semantic architecture that jointly optimizes source coding and
channel adaptation; across varying channel conditions the semantic
approach achieves superior image reconstruction over longer ranges than
traditional compression plus error-correction coding. The principle
transfers directly to NLOS-OCC: a semantic encoder can prioritize
safety-critical vehicular messages or structural-inspection features,
transmitting only the meaning-bearing content that the degraded NLOS
channel can support.

At the architectural level, Bostan \cite{bostan2026crosslayer} surveyed
cross-layer AI-native 6G designs that integrate semantic communication,
ISAC, RIS, and optical wireless communication (OWC) as complementary
programmable physical-layer technologies. The five-domain reference
architecture---physical, control/user-plane, edge-intelligence,
AI-native orchestration, and semantic/task---positions OWC (including
OCC) as a short-range, high-capacity, interference-free physical layer
that semantic encoders can exploit when RF bands are congested. For
NLOS-OCC this means the camera receiver may eventually feed a semantic
decoder that extracts task-level information directly from the striped
image, bypassing the bit-level pipeline entirely.

\subsection{Retroreflective Optical ISAC ($\rp{16}$)}
\label{sec:rp16}

Integrated sensing and communication (ISAC) over the optical band is a
rising 6G theme. Retroreflective optical ISAC (RO-ISAC) pairs an
interrogating transceiver with a passive corner-cube retroreflector
(CCR) that returns light along the incident path, enhancing the
reflected signal strength that is the Achilles' heel of NLOS links. Du
\emph{et al.} \cite{du2025roisac} designed a flexible waveform embedding
an OFDM data signal within a maximum-length-sequence (MLS) sensing
frame, jointly delivering communication and ranging/sensing on a single
optical carrier. The CCR's strong back-reflection relaxes the NLOS power
budget, and the OFDM-MLS embedding avoids the separate-sensor overhead
of conventional ISAC. For NLOS-OCC this is significant: a retroreflector
placed on a vehicle or infrastructure element can return enough light
for a camera to decode both the data and a sensing measurement (range,
velocity) from a single captured frame, merging $\rp{1}$ channel
characterization with $\rp{12}$ ISAC in one physical mechanism.

% =============================================================================
\section{The IEEE 802.15.7a-2024 Standard}
\label{sec:standard}

The IEEE 802.15.7a-2024 amendment~\cite{ieee802157a2024}, approved
2024-12-11 and published 2025-02-14, defines PHY~VII---a dedicated OCC
physical layer with seven modes designed to span the full range of
applications from low-rate IoT to high-rate vehicular communication:

\begin{enumerate}
  \item \textbf{RS-OFDM} --- Reed-Solomon outer code with OFDM, up to
    9.6\,Mbps, the highest-rate mode.
  \item \textbf{MIMO-COOK} --- MIMO with colour-orthogonal OOK, using
    the R/G/B channels as independent streams.
  \item \textbf{O-NOMA} --- Optical NOMA with power-domain
    superposition and SIC, supporting multi-user access on a single
    optical channel.
  \item \textbf{MIMO-OOK} --- Spatial-multiplexing MIMO with per-LED
    OOK, the camera resolving the LEDs by position.
  \item \textbf{HOOK-OFDM} --- Hybrid OOK with OFDM, combining the
    robustness of OOK with the spectral efficiency of OFDM.
  \item \textbf{HS2PSK-OFDM} --- Hybrid S2-PSK (spatial two-phase shift
    keying) with OFDM, a spatial-colour hybrid.
  \item \textbf{BPPM} --- Binary pulse position modulation, the most
    robust mode for low-SNR NLOS channels.
\end{enumerate}

The standard targets a 200\,m communication range and 350\,km/h mobility,
encompassing highway V2V scenarios. Two forward-looking provisions are
notable: an \emph{AI-PHY} framework that allows machine-learning-based
transceiver components (codifying the deep-learning approaches of
$\rp{5}$ within the standard), and a \emph{relaying mechanism} that
enables multi-hop OCC links---the standardized counterpart of the
hybrid VLC+RF fusion in \cite{okasha2022hybrid}. The base
standard~\cite{ieee8021572018} defines PHY~I--III (VLC with
photodiode receivers) and PHY~IV--VI (OCC), against which the amendment
is defined.

% =============================================================================
\section{Open Challenges and a Prioritized Roadmap}
\label{sec:roadmap}

The literature reveals a clear dependency chain:
$\rp{1}\rightarrow\rp{4}/\rp{6}\rightarrow\rp{9}\rightarrow\rp{11}$.
Accurate channel models ($\rp{1}$) define the signal characteristics
that recovery algorithms ($\rp{4}$, $\rp{6}$) must handle; recovery
performance determines whether MIMO and network-level techniques
($\rp{9}$) can operate on the recovered symbols; and the network layer
enables the applications ($\rp{11}$) that justify the effort. We
prioritize the research points accordingly.

\begin{table}[t]
\centering
\caption{Prioritized research roadmap.}
\label{tab:roadmap}
\small
\begin{tabular}{l l l}
\toprule
\textbf{Horizon} & \textbf{RPs} & \textbf{Rationale} \\
\midrule
Near-term & $\rp{4},\rp{7},\rp{12}$ & Denoising + smartphone = deployable demos \\
Mid-term  & $\rp{1},\rp{5},\rp{9},\rp{11}$ & Channel + DL + MIMO $\Rightarrow$ V2X \\
Long-term & $\rp{6},\rp{8},\rp{10}$ & Event cameras, RIS, 6G fusion \\
Frontier  & $\rp{13},\rp{14},\rp{15},\rp{16}$ & UAV, underwater, semantic 6G, RO-ISAC \\
\bottomrule
\end{tabular}
\end{table}

\paragraph{Near-term ($\rp{4}, \rp{7}, \rp{12}$).}
The fastest path to a deployable NLOS-OCC demonstration combines
adaptive thresholding ($\rp{4}$) with a smartphone receiver ($\rp{7}$)
in an indoor IoT or VLP scenario ($\rp{12}$). The required components
exist; the work is integration, robustness testing, and API-level
access to the RAW sensor stream (as demonstrated by
Xiao \emph{et al.}~\cite{winklink2024}, whose two-stage DNN achieves
robust demodulation behind complex, unseen backgrounds).

\paragraph{Mid-term ($\rp{1}, \rp{5}, \rp{9}, \rp{11}$).}
The vehicular use case ($\rp{11}$) demands accurate NLOS channel models
($\rp{1}$), deep-learning recovery robust to motion and weather
($\rp{5}$), and MIMO multiplexing for adequate throughput ($\rp{9}$).
The IEEE 802.15.7a modes provide the interoperability layer; the
research challenge is closing the gap between laboratory demonstrations
and the 200\,m / 350\,km/h targets of the standard.

\paragraph{Long-term ($\rp{6}, \rp{8}, \rp{10}$).}
Event cameras ($\rp{6}$) dissolve the sampling bottleneck but require a
fundamentally new processing stack; RIS ($\rp{8}$) adds coverage but
introduces control-plane complexity; and heterogeneous 6G fusion
($\rp{10}$) requires cross-layer orchestration of OCC, RF, and sensing.
These are the directions with the highest potential payoff and the
highest uncertainty.

\paragraph{Frontier ($\rp{13}, \rp{14}, \rp{15}, \rp{16}$).}
Four emerging directions extend NLOS-OCC beyond its terrestrial,
bit-level origins. UAV-mounted OCC ($\rp{13}$) \cite{chang2026uavocc,
kondo2024farmocc} adds a mobile aerial receiver that can serve
infrastructure-free sensor networks; the challenge is trajectory
optimization and wide-FOV optics. Underwater and cross-medium OCC
($\rp{14}$) \cite{hamidnejad2023uwmimoocc, tsai2024uavwatertoair} must
contend with wavelength-dependent absorption and the turbulent
air-water interface, but enables AUV telemetry and ocean-IoT links.
Semantic optical communication ($\rp{15}$)
\cite{nennouche2025semanticuwoc, bostan2026crosslayer} reframes the
receiver as a meaning extractor rather than a bit decoder, trading
fidelity for task-level robustness---a natural fit for the
bandwidth-limited NLOS channel. Retroreflective optical ISAC
($\rp{16}$) \cite{du2025roisac} merges communication and sensing on a
single optical carrier using passive retroreflectors, directly
addressing the NLOS power-budget deficit. These directions are at the
concept-to-prototype stage; their convergence with the core framework
will define the next generation of NLOS-OCC research.

\paragraph{Cross-cutting challenges.}
Three issues cut across all horizons: (i) \emph{training data} for
deep-learning NLOS-OCC is scarce and expensive to collect, motivating
synthetic data generation via ray tracing and GAN-based augmentation;
(ii) \emph{synchronization} between the LED transmitter and the camera
receiver remains non-trivial, especially under motion, and
LightSync~\cite{lightsync2013}-style ambient-light synchronization
deserves renewed attention; and (iii) \emph{standardization compliance}
---ensuring that research prototypes can interoperate under the
802.15.7a-2024 framework---will determine whether the field's advances
translate into deployed systems.

% =============================================================================
\section{Conclusion}
\label{sec:conclusion}

Non-Line-of-Sight Optical Camera Communication sits at the intersection
of optical wireless, computational imaging, and mobile computing. Its
core challenge---recovering MHz-modulated data from a reflection-challenged,
ambient-light-flooded, 30-fps image stream---has been attacked from
complementary directions: pilot-aided channel modelling and coding at
the physical layer; U-Net and GAN-based deep-learning recovery at the
signal layer; learning-based smartphone demodulation at the hardware layer;
MIMO-NOMA and hybrid RF fusion at the network layer; and event cameras
that rethink the sensing primitive itself. The IEEE 802.15.7a-2024
standard now provides the interoperability framework, and the
dependency chain $\rp{1}\rightarrow\rp{4}/\rp{6}\rightarrow\rp{9}
\rightarrow\rp{11}$ identifies where the next advances must come from.
Four emerging directions---UAV-mounted OCC ($\rp{13}$), underwater and
cross-medium OCC ($\rp{14}$), semantic optical communication toward
AI-native 6G ($\rp{15}$), and retroreflective optical ISAC
($\rp{16}$)---extend the framework to aerial, aquatic, semantic, and
joint-sensing frontiers. The near-term opportunity is the smartphone-based indoor demo; the
mid-term prize is the vehicular V2X link at the standard's 200\,m /
350\,km/h target; and the long-term frontier is the event-camera-equipped,
RIS-assisted, 6G-fused, semantically-aware NLOS-OCC network spanning
air, water, and ground.

% =============================================================================
\bibliographystyle{plain}
\bibliography{refs}

\end{document}
